\documentclass[12pt,a4paper]{article}

\usepackage[utf-8]{inputenc}
\usepackage[portuguese]{babel}
\usepackage[margin=2cm]{geometry}
\usepackage{graphicx}
\usepackage{listings}
\usepackage{xcolor}
\usepackage{hyperref}
\usepackage{amsmath}
\usepackage{amssymb}
\usepackage{fancyhdr}
\usepackage{booktabs}
\usepackage{float}
\usepackage{xspace}

% Configuração de código
\lstset{
    language=C,
    basicstyle=\ttfamily\small,
    keywordstyle=\color{blue},
    commentstyle=\color{gray},
    stringstyle=\color{red},
    breaklines=true,
    tabsize=4,
    showstringspaces=false,
    framexleftmargin=5pt,
    framerule=0.5pt,
    backgroundcolor=\color{lightgray!10},
    numbers=left,
    numberstyle=\tiny\color{gray},
    stepnumber=1,
    xleftmargin=15pt
}

% Estilos de página
\pagestyle{fancy}
\fancyhf{}
\rhead{Testes Unitários em C/C++ para Arduino}
\lhead{Curso Seguidor Spirit}
\rfoot{\thepage}
\lfoot{\today}

\title{
    \Huge \textbf{Testes Unitários em C/C++ para Arduino} \\
    \large Guia Completo para Testes no Projeto Seguidor
}
\author{Bozelli}
\date{\today}

\begin{document}

\maketitle

\tableofcontents
\newpage

% ========== SEÇÃO 1 ==========
\section{Introdução}

Testes unitários são um componente crítico no desenvolvimento de software embarcado.
Especialmente em projetos Arduino, onde o ciclo de feedback é lento (recompilação $\rightarrow$ upload $\rightarrow$ teste físico),
testes automatizados podem economizar horas de desenvolvimento.

Este documento apresenta um guia prático e completo para implementar testes unitários 
em C/C++ para o projeto do robô seguidor. O foco é na simplicidade e aplicabilidade real.

\subsection{O que você aprenderá}

\begin{itemize}
    \item Conceitos fundamentais de testes unitários
    \item Por que testar é importante em sistemas embarcados
    \item Como escolher o framework certo (Unity)
    \item Setup completo do ambiente de testes
    \item Boas práticas para escrita de testes
    \item Integração com CI/CD
\end{itemize}

% ========== SEÇÃO 2 ==========
\section{Conceitos Básicos}

\subsection{Definição}

Um \textbf{teste unitário} é um teste pequeno e focado que verifica se uma função ou módulo 
funciona corretamente de forma isolada.

\begin{lstlisting}[caption=Exemplo básico de teste unitário]
void test_somar_dois_numeros(void) {
    int resultado = somar(2, 3);
    TEST_ASSERT_EQUAL_INT(5, resultado);
}
\end{lstlisting}

\subsection{Características de um bom teste}

Um teste unitário deve ter as seguintes características:

\begin{table}[H]
\centering
\begin{tabular}{|l|p{8cm}|}
\toprule
\textbf{Característica} & \textbf{Descrição} \\
\midrule
Uma responsabilidade & Testa apenas uma coisa \\
Nome descritivo & Deixa claro o que testa \\
Isolado & Não depende de outros testes \\
Rápido & Executa em milissegundos \\
Determinístico & Sempre passa ou falha (não muda) \\
\bottomrule
\end{tabular}
\caption{Características de um bom teste unitário}
\end{table}

\subsection{Padrão Arrange-Act-Assert}

Todo teste bem estruturado segue este padrão de três seções:

\begin{lstlisting}[caption=Padrão AAA - Arrange-Act-Assert]
void test_motor_velocidade_limitada(void) {
    // ARRANGE: Preparar estado
    inicializarMotor();
    
    // ACT: Executar a função
    controlarMotor(500, 500);  // Tenta acima do máximo
    
    // ASSERT: Verificar resultado
    TEST_ASSERT_LESS_THAN_OR_EQUAL(255, getVelocidade());
}
\end{lstlisting}

% ========== SEÇÃO 3 ==========
\section{Por Que Testar?}

\subsection{Problemas sem testes}

Sem testes unitários, você enfrenta vários problemas:

\begin{itemize}
    \item \textbf{Debugging demorado:} Descobrir bugs apenas após carregar no hardware
    \item \textbf{Refatoração arriscada:} Medo de quebrar algo ao modificar código
    \item \textbf{Ciclo lento:} Recompile $\rightarrow$ Upload $\rightarrow$ Teste $\rightarrow$ Repita
    \item \textbf{Documentação ruim:} Incerteza sobre como as funções devem funcionar
\end{itemize}

\subsection{Vantagens com testes}

Com testes unitários implementados:

\begin{itemize}
    \item \textbf{Feedback instantâneo:} Testes rodam em segundos no computador
    \item \textbf{Confiança para refatorar:} Mude código sabendo que testes validam
    \item \textbf{Documentação viva:} Testes mostram exatamente como usar cada função
    \item \textbf{Bugs detectados cedo:} Antes de enviar ao hardware
\end{itemize}

\subsection{Comparação de tempo}

\begin{table}[H]
\centering
\begin{tabular}{|l|c|c|}
\toprule
\textbf{Tarefa} & \textbf{Sem Testes} & \textbf{Com Testes} \\
\midrule
Escrever código & 5 min & 5 min \\
Testar mudanças & 5 min & 0,1 min \\
Debugar & 10 min & 0,5 min \\
Total & \textbf{20 min} & \textbf{5,6 min} \\
\bottomrule
\end{tabular}
\caption{Comparação de tempo: com vs sem testes}
\end{table}

% ========== SEÇÃO 4 ==========
\section{Frameworks de Teste}

\subsection{Comparação de frameworks}

Existem vários frameworks de teste para C/C++. A escolha depende do projeto:

\begin{table}[H]
\centering
\begin{tabular}{|l|c|c|c|c|}
\toprule
\textbf{Framework} & \textbf{Melhor para} & \textbf{Tamanho} & \textbf{Dificuldade} & \textbf{C puro} \\
\midrule
Unity & Embarcados & Pequeno & Fácil & Sim \\
ArduinoUnit & Arduino IDE & Pequeno & Fácil & Sim \\
Google Test & Projetos grandes & Grande & Difícil & Não \\
CppUTest & C/C++ embarcado & Médio & Médio & Sim \\
\bottomrule
\end{tabular}
\caption{Comparação de frameworks de teste}
\end{table}

\subsection{Por que Unity?}

Para o projeto do seguidor, escolhemos \textbf{Unity} pelos seguintes motivos:

\begin{enumerate}
    \item \textbf{Leve:} Apenas $\sim$700 linhas de código
    \item \textbf{Simples:} Aprende em 5 minutos
    \item \textbf{C puro:} Compatível com código Arduino
    \item \textbf{Multiplataforma:} Windows, Linux, Mac
    \item \textbf{Sem dependências:} Funciona em qualquer lugar
\end{enumerate}

% ========== SEÇÃO 5 ==========
\section{Setup do Unity}

\subsection{Passo 1: Baixar Unity}

\begin{lstlisting}[language=bash, caption=Baixar Unity via Git]
cd seguidor-projeto
mkdir -p test
cd test
git clone https://github.com/ThrowTheSwitch/Unity.git
\end{lstlisting}

\subsection{Passo 2: Estrutura de pastas}

O projeto deve ficar organizado assim:

\begin{lstlisting}[language=bash, caption=Estrutura recomendada]
seguidor-projeto/
├── src/
│   ├── motor/
│   │   ├── motor.h
│   │   └── motor.c
│   ├── pid/
│   │   ├── pid.h
│   │   └── pid.c
│   ├── sensores/
│   │   ├── sensores.h
│   │   └── sensores.c
│   └── config/
│       └── config.h
├── test/
│   ├── Unity/          # Framework (git clone)
│   ├── mock_arduino.h  # Simulações
│   ├── mock_arduino.c
│   ├── test_motor.c    # Testes
│   ├── test_pid.c
│   └── run_tests.sh    # Script
└── README.md
\end{lstlisting}

\subsection{Passo 3: Mock do Arduino}

Como os testes rodam no computador (não no Arduino), você precisa simular as funções do Arduino.

\begin{lstlisting}[caption=mock\_arduino.h - Interface dos mocks]
#ifndef MOCK_ARDUINO_H
#define MOCK_ARDUINO_H

#define HIGH 1
#define LOW 0
#define OUTPUT 1
#define INPUT 0

typedef struct {
    int pinMode_calls;
    int digitalWrite_calls;
    int ledcWrite_calls;
} MockState;

extern MockState mock_state;

void resetMocks(void);
void pinMode(int pin, int mode);
void digitalWrite(int pin, int value);
void ledcWrite(int channel, int value);

#endif
\end{lstlisting}

\begin{lstlisting}[caption=mock\_arduino.c - Implementação dos mocks]
#include "mock_arduino.h"

MockState mock_state = {0};

void resetMocks(void) {
    mock_state.pinMode_calls = 0;
    mock_state.digitalWrite_calls = 0;
    mock_state.ledcWrite_calls = 0;
}

void pinMode(int pin, int mode) {
    mock_state.pinMode_calls++;
}

void digitalWrite(int pin, int value) {
    mock_state.digitalWrite_calls++;
}

void ledcWrite(int channel, int value) {
    mock_state.ledcWrite_calls++;
}
\end{lstlisting}

% ========== SEÇÃO 6 ==========
\section{Suite de Testes Prática}

\subsection{Estrutura geral}

\begin{lstlisting}[caption=Estrutura básica de um arquivo de teste]
#include "unity.h"
#include "mock_arduino.h"
#include "../src/motor/motor.h"

void setUp(void) {
    // Executado ANTES de cada teste
    resetMocks();
    inicializarMotor();
}

void tearDown(void) {
    // Executado DEPOIS de cada teste
    pararMotores();
}

void test_motor_inicializa(void) {
    TEST_ASSERT_GREATER_THAN(0, mock_state.pinMode_calls);
}

int main(void) {
    UNITY_BEGIN();
    RUN_TEST(test_motor_inicializa);
    return UNITY_END();
}
\end{lstlisting}

\subsection{Exemplo completo: Testes do Motor}

Vamos testar um módulo de motor com velocidade limitada entre 0-255:

\begin{lstlisting}[caption=test\_motor.c - Suite completa de testes]
#include "unity.h"
#include "mock_arduino.h"
#include "../src/motor/motor.h"

void setUp(void) {
    resetMocks();
    inicializarMotor();
}

void tearDown(void) {
    pararMotores();
}

// Testes de velocidade normal
void test_velocidade_50(void) {
    controlarMotor(50, 50);
    TEST_ASSERT_EQUAL_INT(50, getVelocidadeMotorEsquerdo());
}

// Teste de velocidade acima do máximo
void test_velocidade_acima_maximo(void) {
    controlarMotor(500, 500);
    TEST_ASSERT_LESS_THAN_OR_EQUAL(255, 
        getVelocidadeMotorEsquerdo());
}

// Teste de velocidade negativa
void test_velocidade_negativa(void) {
    controlarMotor(-50, -50);
    TEST_ASSERT_EQUAL_INT(0, getVelocidadeMotorEsquerdo());
}

// Teste de independência dos motores
void test_motores_independentes(void) {
    controlarMotor(100, 200);
    TEST_ASSERT_EQUAL_INT(100, 
        getVelocidadeMotorEsquerdo());
    TEST_ASSERT_EQUAL_INT(200, 
        getVelocidadeMotorDireito());
}

int main(void) {
    UNITY_BEGIN();
    RUN_TEST(test_velocidade_50);
    RUN_TEST(test_velocidade_acima_maximo);
    RUN_TEST(test_velocidade_negativa);
    RUN_TEST(test_motores_independentes);
    return UNITY_END();
}
\end{lstlisting}

\subsection{Resultado esperado}

\begin{lstlisting}[language=bash, caption=Saída dos testes]
test_velocidade_50...PASS
test_velocidade_acima_maximo...PASS
test_velocidade_negativa...PASS
test_motores_independentes...PASS

-----------------------
4 Tests 4 Passed 0 Failed
-----------------------
\end{lstlisting}

% ========== SEÇÃO 7 ==========
\section{Boas Práticas}

\subsection{1. Um conceito por teste}

Cada teste deve verificar apenas uma coisa. Se falhar, você sabe exatamente o quê quebrou.

\begin{lstlisting}[caption=Errado vs Correto]
// ERRADO: Múltiplos conceitos
void test_motor(void) {
    controlarMotor(100, 100);
    TEST_ASSERT_EQUAL_INT(100, getVelocidade());
    pararMotores();
    TEST_ASSERT_EQUAL_INT(0, getVelocidade());
}

// CORRETO: Um conceito por teste
void test_motor_controla_velocidade(void) {
    controlarMotor(100, 100);
    TEST_ASSERT_EQUAL_INT(100, getVelocidade());
}

void test_motor_para_com_sucesso(void) {
    pararMotores();
    TEST_ASSERT_EQUAL_INT(0, getVelocidade());
}
\end{lstlisting}

\subsection{2. Nomes descritivos}

Nomes de testes devem descrever exatamente o que testam.

\begin{lstlisting}[caption=Nomes de testes]
// Ruim
void test1(void) { }
void test_motor(void) { }

// Bom
void test_motor_limita_velocidade_em_255(void) { }
void test_motor_para_ambos_motores_com_sucesso(void) { }
\end{lstlisting}

\subsection{3. Teste casos extremos}

Sempre teste limites: mínimo, máximo, negativo, overflow.

\begin{lstlisting}[caption=Cobertura de casos extremos]
void test_velocidade_minima(void) {
    controlarMotor(0, 0);
    TEST_ASSERT_EQUAL_INT(0, getVelocidade());
}

void test_velocidade_maxima(void) {
    controlarMotor(255, 255);
    TEST_ASSERT_EQUAL_INT(255, getVelocidade());
}

void test_velocidade_overflow(void) {
    controlarMotor(256, 256);
    TEST_ASSERT_LESS_THAN_OR_EQUAL(255, getVelocidade());
}
\end{lstlisting}

\subsection{4. Use setUp() e tearDown()}

Evite código repetido usando estas funções:

\begin{lstlisting}[caption=setUp e tearDown]
void setUp(void) {
    // Executado ANTES de cada teste
    // Use para inicializar estado comum
    resetMocks();
    inicializarMotor();
}

void tearDown(void) {
    // Executado DEPOIS de cada teste
    // Use para limpar recursos
    pararMotores();
}
\end{lstlisting}

% ========== SEÇÃO 8 ==========
\section{Executando Testes}

\subsection{Opção 1: Compilação manual}

\begin{lstlisting}[language=bash, caption=Compilar manualmente]
cd test
gcc -I. -I../src -DTEST \
    test_motor.c \
    ../src/motor/motor.c \
    mock_arduino.c \
    Unity/src/unity.c \
    -o test_motor

./test_motor
\end{lstlisting}

\subsection{Opção 2: Script bash}

\begin{lstlisting}[language=bash, caption=run\_tests.sh - Automatizar testes]
#!/bin/bash

DIR="$(cd "$(dirname "$0")" && pwd)"
cd "$DIR"

echo "Compilando testes..."
gcc -I. -I../src -DTEST \
    test_motor.c \
    ../src/motor/motor.c \
    mock_arduino.c \
    Unity/src/unity.c \
    -o test_motor

if [ $? -ne 0 ]; then
    echo "Erro ao compilar"
    exit 1
fi

echo "Rodando testes..."
./test_motor
\end{lstlisting}

Use com:
\begin{lstlisting}[language=bash]
chmod +x test/run_tests.sh
./test/run_tests.sh
\end{lstlisting}

\subsection{Opção 3: CMake}

Para projetos maiores, use CMake:

\begin{lstlisting}[caption=CMakeLists.txt]
cmake_minimum_required(VERSION 3.10)
project(SeguidorTests)

add_library(unity STATIC Unity/src/unity.c)
target_include_directories(unity PUBLIC Unity/src)

function(add_test_suite name file)
    add_executable(${name} ${file} 
        ../src/motor/motor.c
        mock_arduino.c
    )
    target_include_directories(${name} PRIVATE . ../src)
    target_link_libraries(${name} unity)
    add_test(NAME ${name} COMMAND ${name})
endfunction()

add_test_suite(test_motor test_motor.c)
enable_testing()
\end{lstlisting}

Use com:
\begin{lstlisting}[language=bash]
cd test/build
cmake ..
make
make test
\end{lstlisting}

% ========== SEÇÃO 9 ==========
\section{Integração com Git e CI/CD}

\subsection{GitHub Actions}

Configure testes automáticos a cada push:

\begin{lstlisting}[language=yaml, caption=.github/workflows/tests.yml]
name: Unit Tests

on: [push, pull_request]

jobs:
  test:
    runs-on: ubuntu-latest
    
    steps:
    - uses: actions/checkout@v2
    
    - name: Run tests
      run: |
        cd test
        mkdir -p build
        cd build
        cmake ..
        make
        make test
\end{lstlisting}

Agora, a cada push, os testes rodam automaticamente.

\subsection{Checklist pré-commit}

\begin{lstlisting}[language=bash, caption=.git/hooks/pre-commit]
#!/bin/bash
echo "Rodando testes..."
./test/run_tests.sh

if [ $? -ne 0 ]; then
    echo "Testes falharam! Commit cancelado."
    exit 1
fi
\end{lstlisting}

% ========== SEÇÃO 10 ==========
\section{Resumo}

\subsection{Checklist para implementar testes}

\begin{enumerate}
    \item[\checkmark] Baixe e configure Unity
    \item[\checkmark] Estruture pastas (src/, test/)
    \item[\checkmark] Crie mock\_arduino.h e mock\_arduino.c
    \item[\checkmark] Escreva testes para cada módulo
    \item[\checkmark] Siga padrão AAA (Arrange-Act-Assert)
    \item[\checkmark] Use nomes descritivos
    \item[\checkmark] Teste casos extremos
    \item[\checkmark] Configure script de testes
    \item[\checkmark] Integre com git (hooks ou CI/CD)
\end{enumerate}

\subsection{Próximos passos}

\begin{itemize}
    \item Teste cada módulo: motor, sensores, PID, Bluetooth
    \item Crie mocks mais avançados conforme necessário
    \item Configure CI/CD para rodar testes automaticamente
    \item Documente padrões de teste no projeto
\end{itemize}

\subsection{Referências}

\begin{itemize}
    \item Unity Documentation: \url{https://github.com/ThrowTheSwitch/Unity}
    \item Test-Driven Development for Embedded C (Pragmatic Programmers)
    \item Arduino Best Practices
\end{itemize}

\end{document}
